\section{Threat Model and Security Boundaries}\label{sec:threats}

\subsection{Three different kinds of protection}
It is easy to read ``encrypted files in private repositories'' as one guarantee. It is three, with
different owners, and GhostVault provides only the first.

\emph{The content of a file is protected by cryptography.} Its confidentiality and integrity hold
against anyone who can read or change the repositories, GitHub included, as long as the keys stay
secret and AES-GCM is secure. Nothing about this depends on the repositories being private.

\emph{Who can read the repositories is decided by GitHub's access control.} ``Private'' is a
GitHub setting, not encryption. GhostVault refuses a repository that is public, archived or not
writable with the given token when it is added (\secref{sec:gh-checks}), but afterwards GitHub
alone decides who can read it: collaborators, holders of any token with access, GitHub itself,
and everyone if the repository is ever made public. Because of the encryption, what such readers
gain is limited to the metadata in \tabref{tab:learn}; but they do gain that.

\emph{Whether the data is still there is decided by GitHub.} Each chunk is stored exactly once.
If a repository is deleted, an account suspended, a token's access withdrawn for good or the data
lost, every ghost with a chunk there is lost; GitHub's acceptable use policies reserve the right
to ``suspend your Account, throttle your file hosting, or otherwise limit your activity'' when
bandwidth use is ``significantly excessive''~\cite{gh-aup}. GhostVault detects missing or damaged
data and refuses to restore it wrongly; it cannot bring it back. And once a file is ghosted, the
ghost is its only copy unless the user keeps others.

\begin{table*}[t]
\caption{Who the design considers, what they can do, and what stops them}\label{tab:parties}
\centering\footnotesize
\begin{tabularx}{\textwidth}{@{}L{36mm}L{48mm}Y@{}}
\toprule
\textbf{Party or event} & \textbf{Can} & \textbf{What GhostVault does about it} \\
\midrule
GitHub, or anyone who can read the repositories & read every blob, tree and commit, with their
times and history & only ciphertext under keys they never see; random paths; neutral commit
messages; padding hides exact sizes. They learn what \tabref{tab:learn} lists. \\ \addlinespace[1.6pt]
Someone on the network & read or change traffic & TLS, with GitHub's certificate verified against
the operating system's trust store. \\ \addlinespace[1.6pt]
Someone who changes the data on GitHub & change, swap, truncate or delete chunks & detected: every
chunk is checked against its content id, its SHA-256 and its GCM tag, whose associated data binds
file, position, count, size and padding. A restore fails without writing anything under the file's
name. Deleted data cannot be brought back. \\ \addlinespace[1.6pt]
Crashes, power loss, a full disk, network failures & interrupt any operation at any point & every
step that changes something outside the data folder is journaled first; the next start finishes or
undoes it (\secref{sec:journal}). \\ \addlinespace[1.6pt]
Another program working on the same files & replace, edit or rename the original or the restored
file's name while GhostVault works & the original is moved to a hidden name and checked there before
it is deleted; a restored file is published without replacing anything (Sections~\ref{sec:remove}
and~\ref{sec:publish}). \\ \addlinespace[1.6pt]
Another account on the computer, a stolen disk, a copy of the data folder & read what file
permissions let them read & partly: the data folder is readable by its owner only, but the keys
in it are not encrypted (\secref{sec:keys}); disk encryption and protected backups of the folder
are the user's part. \\ \addlinespace[1.6pt]
Malware running as the user & everything the user can do & not covered: it can read the data
folder and the keychain and watch what is opened. \\ \addlinespace[1.6pt]
A compromised dependency, build tool or release & run code inside GhostVault & partly:
dependencies are audited in CI, build tools are pinned, releases carry SHA-256 checksums, and builds
are signed when a certificate is configured (\secref{sec:builds}). \\
\bottomrule
\end{tabularx}
\end{table*}

\subsection{Parties and events}
\tabref{tab:parties} lists who the design considers and what stops each of them. GhostVault
trusts the operating system, the user's account, the Python runtime, the \code{cryptography}
library and the OpenSSL inside it, and the operating system's certificate store for TLS. It does
not trust GitHub's replies, the state of a file between two system calls, or its own state files
after a crash; each is checked before anything depends on it.

\subsection{What GitHub learns}\label{sec:learn}
\tabref{tab:learn} lists what anyone who can read the repositories learns, and what stays hidden.
Commit messages are always \code{Update}, and paths are random, so neither the tree nor the history
names a file. A network observer sees only that the computer talks to \code{api.github.com}, how
much and when.

\begin{table}[t]
\caption{What readers of the repositories learn}\label{tab:learn}
\centering\footnotesize
\begin{tabular}{@{}L{40mm}L{40mm}@{}}
\toprule
\textbf{They learn} & \textbf{They do not learn} \\
\midrule
that the account stores GhostVault-shaped encrypted data, and how much & file names, types or
content \\ \addlinespace[1.6pt]
each ghost's chunk count, so its size to within 512\,KiB with padding (the default), or exactly
without & the exact size, when padding is on \\ \addlinespace[1.6pt]
when each ghost was uploaded and removed, where, and which chunks were committed together &
original paths, timestamps, or hashes of the plaintext \\ \addlinespace[1.6pt]
whether a ghost has a preview (all previews have one size) & what a preview shows, or its
dimensions \\ \addlinespace[1.6pt]
the account's addresses and active times & keys, nonces or tokens \\
\bottomrule
\end{tabular}
\end{table}
